\section{ComprasNet Cross-Jurisdiction Replication}
\label{app:comprasnet_submission}

This appendix tests the central claims of the paper on a second
procurement system. The main text reports results on BEC--SP
(2009--2019), one state-level platform under the same federal
legal frame. Here we re-run the headline analyses on the federal
ComprasNet panel for 2013--2019 and report what survives, what
attenuates, and what fails to generalise. We do not adjust the
framing of the main text on the basis of the federal results;
this appendix is added so that a reader can locate the source of
any disagreement between the BEC headline numbers and an
alternative panel.

\subsection{Panel construction}
\label{app:comprasnet_panel}

We use the bulk dumps published by the Comptroller-General's
Office (CGU) at \texttt{dadosabertos-download.cgu.gov.br} under
\texttt{Portal\-Da\-Transparencia/saida/licitacoes/\{YYYYMM\}\_Licitacoes.zip}.
The series begins in 2013-01, so the federal-overlap window with
the BEC panel is the seven calendar years 2013--2019. We download
\valFedZipMonths{} monthly archives (\valFedAcquireMiBcompressed{}
MiB compressed), each containing a \texttt{ParticipantesLicitação}
CSV with one row per firm-item-tender participation and a
\texttt{Flag\,Vencedor} (SIM/NÃO) field. We restrict to modalities
5 (Pregão) and 9999 (Pregão SRP); other modalities either have no
real loser tracking in this source or do not exist in the federal
panel (Convite federal is effectively extinct, with on the order
of five participations per month nationwide). The build is
DuckDB-native end to end and produces five canonical parquets in
\texttt{data/processed\_comprasnet/} that mirror the BEC schema
used throughout the main text.

The resulting panel has \valFedPanelRows{} participation rows,
\valFedFirms{} distinct CNPJ-14 firms, and \valFedAlwaysLosers{}
firms with zero wins (\valFedALshare{} per cent of the universe).
Pregão SRP accounts for \valFedSharePregaoSRP{} per cent of
participation rows; the residual is Pregão. The overall winner
share is \valFedWinRateGlobal{} per cent. The IQR rule
$\textit{median} + 1.5 \cdot \textit{IQR}$ applied to the
tenders-count distribution among always-losers gives a federal
threshold of \valFedFLthreshold{} (against \valBECalShare{} for
BEC), and the FL count above that threshold is
\valFedFLcount{} firms (\valFedFLshare{} per cent of
always-losers).

\subsection{CADE ground truth on the federal panel}
\label{app:comprasnet_linkage}

CADE prosecutes cartels across both state and federal procurement.
The CADE-cartels CSV used by the main text contains
\valFedCadeRows{} rows across 12 unique processes, of which
\valFedCadeWithCNPJ{} arrive with CNPJs in the public release.
For the federal panel we enrich an additional \valFedCadeEnrichedBEC{}
firms by fuzzy-matching CADE \emph{razão social} against the
BEC-side crossmatch CSV that already linked CADE firms to CNPJs
through the BEC pipeline. We raise the fuzzy threshold to
\valFedFuzzyThreshold{} (Jaro-Winkler) after detecting one false
positive at the previous default of 0.85, where two firms in the
same sacos-de-lixo cartel were collapsed onto a single CNPJ. We
also enforce a CNPJ-uniqueness rule that demotes any
multi-assigned match to unmatched. After enrichment we have
\valFedCadeEnrichedTotal{} CADE rows with CNPJ. Of those,
\valFedRaizesMatched{} CADE 8-digit root CNPJs are observed in
the federal participation panel, producing \valFedDirects{}
direct-defendant establishment rows (some root CNPJs map to more
than one establishment in the federal data).

The CADE-anchored set of $(\textit{tender}, \textit{item})$ pairs
in which at least one direct defendant participates has
\valFedAnchoredTenders{} elements. The cobidder set obtained by
expanding that anchored set is \valFedCobidders{} federal firms,
which is about an order of magnitude larger than the BEC cobidder
set (193) even though the federal direct-defendant set is
substantially smaller. The reason is mechanical: the cobidder set
scales with the number of anchored tender-items, not with the
number of direct defendants, and federal Pregão tenders typically
have many more bidders per item than BEC tenders. Among federal
cobidders, \valFedALcobidders{} are always-losers and
\valFedFLcobidders{} satisfy the FL$\geq$\valFedFLthreshold{}
threshold. We carry both \emph{case\_scope = national} (the two
roots associated with nationally adjudicated cartels) and
\emph{case\_scope = state} (the twelve roots adjudicated against
state procurement but observed federally) into all downstream
analyses; the headline tables pool both, and a state-only sub-row
is reported as a robustness check where space permits.

The linkage pipeline is intentionally conservative and proceeds in
three audited passes. The RAIS pass (2015--2017 employment-link
records) recovers \valFedCadeEnrichedRAIS{} additional roots out of
the fifteen unmatched after BEC reuse, validated by UF and
razão-social agreement. A second pass against the Receita Federal
Empresas open dump (May 2026 snapshot) recovers
\valFedCadeEnrichedRF{} additional roots, including one identified
through CADE acórdão search (Frontal Indústria e Comércio de
Móveis Hospitalares Ltda, CNPJ raiz 01140694, one of the two firms
of the Frontal Group condemned in case 08012.009732/2008-01 for the
Sanguessuga / Mobile Health Units cartel). Final dictionary
coverage is \valFedCadeEnrichedTotal{} of \valFedCadeRows{} CADE
rows (\valFedCadeCoverage{} per cent). The remaining unmatched
include five summary rows that have no firm-level razão social, and
three firms for which neither RAIS, RF Empresas, nor CADE press
material allow disambiguation (Astéria, Matrix Artefatos Plásticos,
and the second Frontal Group firm). A structural finding from this
exercise: \valFedCadeEnrichedRF{} of the firms added at the RF pass
have zero participation in the federal ComprasNet panel. CADE
prosecutions in this sample are predominantly against
sub-national (state or municipal) procurement cartels; only the
medicamentos-genéricos and unidades-móveis-de-saúde processes
have nominally national scope, and only a subset of those
defendants are observed in federal procurement. The cobidder set
expansion saturates at the BEC-reuse-plus-RAIS step. The full
pipeline is in
\texttt{scripts/65\_cade\_comprasnet\_linkage.py} (BEC reuse),
\texttt{scripts/67\_rais\_enrich.py} (RAIS audit),
\texttt{scripts/68\_promote\_rais\_to\_v2.py} (manual RAIS
promotion),
\texttt{scripts/70\_rf\_enrich.py} (RF Empresas audit),
\texttt{scripts/71\_promote\_rf\_to\_v3.py} (manual RF promotion
including the CADE acórdão match for Frontal Móveis Hospitalares),
and \texttt{scripts/72\_linkage\_v3.py} (re-run linkage on the
enriched set). The audit logs are in
\texttt{data/processed\_comprasnet/cade\_link\_v1/} and
\texttt{cade\_link\_v2/}.
The v3 parquet outputs are in
\texttt{data/processed\_comprasnet/cade\_link\_v3/}.

\subsection{Headline replication}
\label{app:comprasnet_replication}

\paragraph{AN-001 \textemdash{} loser-side concentration.}
The structural sanity check survives. Always-losers account for
\valFedALshare{} per cent of the federal universe, against
\valBECalShare{} per cent for BEC. The share of always-losers
that meet the FL threshold is \valFedFLshare{} per cent federal
against \valBECflShare{} per cent BEC. The federal distribution
of tenders-count among always-losers is more right-skewed than
the BEC counterpart: the p99 is \valFedTcPNinetyNine{} federal against
\valBECtcPNinetyNine{} BEC, the p95 is \valFedTcPNinetyFive{} against 38, and the
maximum is \valFedTcMax{} against 989. Loser-side concentration
is a meaningful concept on this panel, with a thicker upper tail
than on BEC.

\paragraph{AN-004 \textemdash{} cobidder AUC.}
The firm-level AUC of the FL binary classifier against the
always-loser cobidder pool is $\valFedAucBinary{}$ federal
(95\% CI $[\valFedAucBinaryLow{}, \valFedAucBinaryHi{}]$),
against the BEC reproduction of \valBECaucBinary{}
(95\% CI $[\valBECaucBinaryLow{}, \valBECaucBinaryHi{}]$). The
continuous loss-intensity classifier
$\log(1 + \mathrm{tenders\,count})$ improves on the binary in
both panels: federal $\valFedAucContinuous{}$
($[\valFedAucContinuousLow{}, \valFedAucContinuousHi{}]$) versus
binary $\valFedAucBinary{}$, with DeLong $Z =$ \valFedDelongZraw,
$p =$ \valFedDelongPraw. The DeLong direction matches BEC
($p =$ \valBECdelongPraw). The federal absolute AUC is well above
random and the continuous-beats-binary insight of the main text
generalises, but the federal absolute level is markedly below the
BEC level. The gap, about \valFedAucDrop{} in AUC binary, is the
single largest difference between the two panels.

\paragraph{AN-007 \textemdash{} boundary against direct
defendants.}
The structural boundary survives. The firm-level AUC of the
binary FL classifier against direct CADE defendants in the
universe is \valFedDirectsBinary{} federal
($[\valFedDirectsBinaryLow{}, \valFedDirectsBinaryHi{}]$) against
\valBECdirectsBinary{} BEC. The continuous classifier gives
\valFedDirectsContinuous{} federal against \valBECdirectsContinuous{}
BEC; both are below 0.5, consistent with the design statement in
the main text that loser-side ranks do not rank winners and that
direct defendants in CADE adjudications are mostly winners (in
the federal set, \valFedDirectsWinshare{} of \valFedDirects{}
direct-defendant rows are non-AL). The pattern is qualitatively
identical to BEC.

\paragraph{AN-006 \textemdash{} strict prospective holdout.}
We split the federal panel at 2016-12-31 (\valFedTrainRows{} M
rows pre-cut versus \valFedTestRows{} M rows post-cut) and build
the FL classifier using only the train-period activity of each
firm. The train-period IQR threshold is \valFedTrainThreshold{},
the train-period always-loser pool is \valFedTrainAL{} firms, and
the train-period FL classifier is fixed before any test data is
inspected. With cobidder labels defined as in the main text (any
always-loser that ever co-bid with a CADE direct defendant), the
firm-level AUC of the train-period classifier is
\valFedHoldoutBecStyleBinary{} binary
($[\valFedHoldoutBecStyleBinaryLow{}, \valFedHoldoutBecStyleBinaryHi{}]$)
and \valFedHoldoutBecStyleContinuous{} continuous
($[\valFedHoldoutBecStyleContinuousLow{}, \valFedHoldoutBecStyleContinuousHi{}]$).
The BEC analogue reported in the main text is in the
$[\valBECholdoutLow{}, \valBECholdoutHi{}]$ range. The federal
holdout is below the BEC range. We also run an ultra-strict
variant in which the cobidder labels themselves are restricted to
the 2017--2019 test years (this requires not only that the
classifier be temporally honest but that the cobidder behaviour
be observed in the test window). Under this stricter design the
federal AUC drops to \valFedHoldoutUltraStrictBinary{} binary and
\valFedHoldoutUltraStrictContinuous{} continuous. We report both
numbers; the BEC-style number is the direct comparator to the
main text and the ultra-strict number is the deployment-realistic
floor for an agency that would have to re-detect cobidder
behaviour in unseen years.

\paragraph{AN-014 \textemdash{} leakage drop chain.}
We compute three measurements: the in-sample firm-level AUC
($M_1$), the mean AUC across a random 5-fold firm-level partition
($M_2$), and the temporal-holdout AUC reported in the AN-006
paragraph ($M_3$, BEC-style ground truth). The federal values are
$M_1 = \valFedLeakageMOneBinary$ binary and
$\valFedLeakageMOneContinuous$ continuous; $M_2 =
\valFedLeakageMTwoBinary$ binary (SD \valFedLeakageMTwoBinarySD{}
across folds) and $\valFedLeakageMTwoContinuous$ continuous (SD
\valFedLeakageMTwoContinuousSD{}); and $M_3 =
\valFedLeakageMThreeBinary$ binary and $\valFedLeakageMThreeContinuous$
continuous. The $M_1 \to M_2$ drop is essentially zero in both
panels because the FL classifier has no learned parameters and a
random firm-level partition cannot create out-of-sample drift in
that classifier. The $M_2 \to M_3$ drop is
\valFedLeakageDropTwoThreeBinary{} binary and
\valFedLeakageDropTwoThreeContinuous{} continuous, comparable in
magnitude to the BEC in-sample-to-holdout drop reported in the
main text. The drop pattern generalises; the absolute level does
not.

\subsection{Selection and within-cell mechanism decompositions}
\label{app:comprasnet_decomp}

The main text decomposes the positive baseline price coefficient
into a selection channel (cells with higher FL share have
structurally more expensive non-treated items) and a within-cell
mechanism channel (within a cell, FL presence is associated with
lower winner price relative to ceiling, mediated by bidder-count
inflation). We attempt the same decomposition on the federal
panel.

\paragraph{AN-039 \textemdash{} selection mechanism.}
We build a federal item-level panel of \valFedItemPanelRows{}
rows by extracting \texttt{ItemLicita\c{c}\~ao.csv} from the
\valFedZipMonths{} monthly bulk archives, joining to the
participation panel for $n_{firms}$ per item, and tagging
\texttt{is\_treated\_item} as one if any federal FL firm
(\valFedFLcount{} firms with $\geq$ \valFedFLthreshold{}
tenders-count and zero wins) participates in the item. Cells are
defined as (procuring unit code $\times$ item-stem first 12
characters $\times$ year $\times$ modality). The cell FL share is
the share of treated items inside the cell. Among
\valFedItemPanelNonTreated{} non-treated items, OLS of
$\log(\text{valor item})$ on cell FL share with full FE on the
four cell margins gives a coefficient of
$\valFedSelectionCoefFullFE$ (SE $\valFedSelectionSEFullFE$,
$p < 10^{\valFedSelectionLogP}$). Non-treated mean log price
rises monotonically from the zero-share band to the
highest-share band. The selection direction replicates BEC
($\valBECSelectionCoefFullFE$ in the main text), with smaller
magnitude, in line with the more compressed federal price
distribution. We read this as a federal confirmation that
cell-level positive baselines reflect selection of CADE
defendants into structurally expensive cells, not within-cell
mechanism.

\paragraph{AN-040 \textemdash{} within-cell mechanism (does not
replicate).}
The main-text within-cell mechanism test uses
\texttt{winner\_vs\_ref} (the log of winner price over reference
or ceiling price) as dependent variable, isolating the
deviation of the winning bid from the procuring agency's
internal cap. Federal Portal CGU does not publish reference or
ceiling prices, so we substitute the only price object available:
$\log(\text{valor item})$ for the awarded value and an unit-price
variant $\log(\text{valor item}/\text{quantidade item})$. The
coefficient on \texttt{losers} within cells flips sign relative
to BEC: $\valFedMechRawCoef$ in the federal raw-price regression
versus $\valBECMechRawCoef$ in the BEC reference-normalised
regression, and $\valFedMechUnitCoef$ in the federal unit-price
variant. Controlling for $\log(n_{firms})$ does not undo the
positive sign in federal as it does in BEC. The bidder-count
band split is also different: federal coefficients are
positive in all four bands and \emph{grow} from sparse
(1--3 bidders, $\valFedMechBandSparseCoef$) to dense
(11+ bidders, $\valFedMechBandDenseCoef$), whereas the BEC
counterparts flip from positive in sparse
($\valBECMechBandSparseCoef$) to negative in dense
($\valBECMechBandDenseCoef$). The one BEC primitive that does
replicate is the M1 leg: $\log(n_{firms}) \sim \text{losers}$
gives $\valFedMechNfirmsOnLosersCoef$ federal against
$\valBECMechNfirmsOnLosersCoef$ BEC, confirming that FL
presence inflates the bidder count by roughly two-thirds of a
log-unit in both panels.

We read the AN-040 federal failure as primarily a measurement
gap rather than a substantive contradiction. The main-text
mechanism interpretation is identified off the
winner-versus-reference object; without a reference price the
federal regression conflates within-cell selection
(more-expensive items inside a cell attract FL participation)
with the within-cell mechanism (cover bidding deflating the
winning bid relative to the ceiling). The BEC test discriminates
the two because the ceiling price absorbs the level of the cell
mean; the federal test cannot. Reporting both speccs, the
selection result generalises (AN-039) and the within-cell
mechanism result does not (AN-040 raw); but the bidder-count
inflation leg of the BEC mechanism story (M1) does replicate.

\subsection{Threshold sensitivity}
\label{app:comprasnet_threshold}

The federal IQR rule selects threshold \valFedFLthreshold{}, but
that is not where the federal binary AUC peaks. A sweep over
$\mathrm{tenders\,count} \geq T$ for $T \in \{5, 9, 14, 15, 20,
28, 32, 40, 62, 78, 103, 120, 145, 182, 239, 390\}$ peaks at
$T = \valFedPeakThreshold{}$ with binary AUC
\valFedPeakAUC{}. Above the peak the AUC declines monotonically
to random levels by $T = 390$. The continuous classifier
($\log(1 + \mathrm{tenders\,count})$, threshold-free) reaches
\valFedAucContinuous{} on the same federal pool, dominating any
binary cut. Two implications follow. First, the federal
operational rule should not be the IQR rule transported from BEC;
the federal IQR rule underweights the right tail relative to
where the cobidder concentration peaks. Second, the
continuous-beats-binary insight that the main text develops on
BEC is more pronounced federally: the binary classifier loses
about five percentage points of AUC against the continuous
classifier at the federal peak threshold, against about one to
two percentage points on BEC.

\subsection{Honest interpretation}
\label{app:comprasnet_interpretation}

The federal panel preserves the qualitative structure of the
main-text results and degrades the absolute magnitudes by
something on the order of $\valFedAucDrop$ in AUC binary. We
state the four claims that survive and the three that do not.

\paragraph{Survives.} The loser-side concentration concept is
panel-agnostic (AN-001). The FL classifier discriminates
cobidders from non-cobidders well above chance and with tight
confidence intervals (AN-004). The continuous loss-intensity
classifier dominates the binary classifier on both panels, with
the DeLong direction preserved. The structural boundary against
direct defendants is preserved (AN-007). The drop pattern from
in-sample to temporal holdout is similar across panels
(AN-006, AN-014). The cell-level selection mechanism is
positive and significant (AN-039). The bidder-count inflation
leg of the within-cell mechanism (M1) replicates: FL presence
adds about two-thirds of a log-unit to the number of bidders
inside the cell.

\paragraph{Does not survive.} The AUC level of about $0.92$
reported in the main text is panel-specific; the federal binary
in-sample number is \valFedAucBinary{} and the federal BEC-style
temporal-holdout number is \valFedHoldoutBecStyleBinary{}. The
IQR rule applied to always-loser tenders-count is not a universal
threshold rule; the federal peak is at $T = \valFedPeakThreshold{}$
rather than at the IQR-derived $T = \valFedFLthreshold{}$. The
screen does not transfer plug-and-play across procurement
panels. The within-cell winner-versus-reference mechanism
(AN-040) does not replicate using $\log(\text{valor item})$ as
a substitute dependent variable; the federal data lake does not
expose the reference/ceiling price required to identify the
mechanism cleanly. The bidder-count inflation channel (M1) is
the only AN-040 leg that generalises.

\paragraph{Framing.}
We position the federal numbers as a deployment-realistic floor
rather than as evidence that the main-text claim is wrong. The
BEC in-sample AUC is a ceiling under a single panel with a
particular institutional design. The BEC holdout AUC of about
$0.79$--$0.85$ is closer to the level an agency could expect on
unseen years \emph{within the same panel}. The federal AUC of
about $0.60$ is closer to what an agency could expect on a
different panel under temporal discipline. The widening of the
gap with each level of out-of-sample discipline is the natural
behaviour of a screen whose theoretical content is robust but
whose calibration is institution-coupled.

The paper does not modify its central claims on the basis of
this appendix. The main-text construct definitions, identification
logic, and welfare interpretation are written for BEC and refer
explicitly to BEC sample boundaries. The contribution of this
appendix is to give a reader a falsifiable second-panel
benchmark, to expose the screen to a panel where it could have
failed entirely, and to document that the qualitative theory
survives the test while the calibration constants do not.

\subsection{Replication}
\label{app:comprasnet_replication_artifacts}

All federal numbers in this appendix are regenerable from
artefacts under \texttt{data/processed\_comprasnet/} and
\texttt{work/v20-comprasnet/output/}. The build script is
\texttt{scripts/00\_build\_eventlevel\_comprasnet.py}; the
year-stamped panel for temporal splits is built by
\texttt{scripts/66\_year\_stamped\_panel.py}; the CADE linkage is
in \texttt{scripts/65\_cade\_comprasnet\_linkage.py}; and the
five AN replications are in
\texttt{work/v20-comprasnet/scripts/an\{001,004,004b,006,007,014\}\_*.R}.
Each script writes JSON and CSV outputs in a per-AN subdirectory.
A consolidated comparison table is in
\texttt{work/v20-comprasnet/output/CONSOLIDATED\_RESULTS.md}.
